%% =============================================================
%%  chapters/08_results.tex
%% =============================================================
\chapter{Results and Evaluation}
\label{ch:results}

\section{Experimental Setup and Evaluation Protocol}
\label{sec:experimental-setup}
All results reported in this chapter are computed on the held-out test
split of the dataset described in Section~\ref{sec:final-dataset-composition},
using the model checkpoint selected by best validation performance as
described in Section~\ref{sec:training-configuration}. Occupancy and
passenger classification are evaluated using accuracy and F1-score,
weight estimation is evaluated using mean absolute error and weighted
root mean squared error (WRMSE), and pose estimation is evaluated using
mean per-keypoint error or classification accuracy depending on the final
pose representation chosen in Section~\ref{sec:pose-head}.
%% TODO: Confirm final evaluation metrics per task and any statistical
%% significance testing procedure used.

\section{Occupancy Detection Results}
\label{sec:occupancy-results}
The fused model is expected to achieve high occupancy detection accuracy
given the strong, complementary occupancy signal present in both the
camera and pressure modalities. Table~\ref{tab:occupancy-results}
reports placeholder occupancy detection results across the evaluated
models.

\begin{table}[H]
\centering
\caption{Occupancy detection results on the held-out test set.}
\label{tab:occupancy-results}
\begin{tabular}{@{}lcc@{}}
\toprule
\textbf{Model} & \textbf{Accuracy (\%)} & \textbf{F1-score} \\
\midrule
Camera-only   & \todo{XX.X} & \todo{0.XX} \\
Pressure-only & \todo{XX.X} & \todo{0.XX} \\
Fused (ours)  & \todo{XX.X} & \todo{0.XX} \\
\bottomrule
\end{tabular}
\end{table}
%% TODO: Insert actual occupancy detection results once experiments are run.

\section{Passenger Classification Results}
\label{sec:classification-results}
Passenger classification is expected to benefit most from fusion in cases
where an object placed on the seat produces a pressure signature
ambiguous with a small child but is visually distinguishable from the
camera view, or vice versa under partial visual occlusion. Figure
\ref{fig:classification-confusion} presents a placeholder confusion
matrix for the fused model across all passenger classes.

\begin{figure}[H]
\centering
%% TODO: Replace with the actual confusion matrix figure once generated.
\includegraphics[width=0.7\textwidth]{classification_confusion_matrix}
\caption{Confusion matrix for passenger classification using the fused
camera-pressure model on the held-out test set.}
\label{fig:classification-confusion}
\end{figure}
%% TODO: Discuss the most frequent misclassification pairs once available.

\section{Weight Estimation Results}
\label{sec:weight-results}
Table~\ref{tab:weight-results} reports placeholder weight estimation
error metrics, comparing the fused model against single-modality
baselines.

\begin{table}[H]
\centering
\caption{Weight estimation error on the held-out test set.}
\label{tab:weight-results}
\begin{tabular}{@{}lcc@{}}
\toprule
\textbf{Model} & \textbf{MAE (kg)} & \textbf{WRMSE (kg)} \\
\midrule
Camera-only   & \todo{X.X} & \todo{X.X} \\
Pressure-only & \todo{X.X} & \todo{X.X} \\
Fused (ours)  & \todo{X.X} & \todo{X.X} \\
\bottomrule
\end{tabular}
\end{table}
%% TODO: Insert actual weight estimation results once experiments are run.

\section{Pose Estimation Results}
\label{sec:pose-results}
Pose estimation performance is reported in terms of the metric selected
in Section~\ref{sec:pose-head} (mean per-keypoint error or classification
accuracy). Figure~\ref{fig:pose-qualitative} shows placeholder qualitative
examples comparing predicted and ground-truth pose across representative
seating configurations.

\begin{figure}[H]
\centering
%% TODO: Replace with actual qualitative pose estimation result figures.
\includegraphics[width=0.85\textwidth]{pose_qualitative_results}
\caption{Qualitative comparison of predicted and ground-truth occupant
pose for representative seating configurations.}
\label{fig:pose-qualitative}
\end{figure}
%% TODO: Insert quantitative pose estimation results table once available.

\section{Ablation Study: Modality Contribution}
\label{sec:ablation-study}
To quantify the individual contribution of each modality, the fused model
is compared against camera-only and pressure-only variants trained with
the same task heads and training configuration. Table
\ref{tab:ablation-modality} summarises this comparison across all four
output tasks.

\begin{table}[H]
\centering
\caption{Ablation study comparing camera-only, pressure-only, and fused
model variants across all output tasks.}
\label{tab:ablation-modality}
\begin{tabular}{@{}lccc@{}}
\toprule
\textbf{Metric} & \textbf{Camera-only} & \textbf{Pressure-only} & \textbf{Fused} \\
\midrule
Occupancy accuracy (\%)      & \todo{XX.X} & \todo{XX.X} & \todo{XX.X} \\
Classification accuracy (\%) & \todo{XX.X} & \todo{XX.X} & \todo{XX.X} \\
Weight MAE (kg)               & \todo{X.X}  & \todo{X.X}  & \todo{X.X} \\
Pose error                    & \todo{X.X}  & \todo{X.X}  & \todo{X.X} \\
\bottomrule
\end{tabular}
\end{table}
%% TODO: Insert actual ablation results once experiments are complete, and
%% discuss statistical significance of the fused model's improvement.

\section{Sim2Real Transfer Evaluation}
\label{sec:sim2real-transfer-evaluation}
To evaluate Sim2Real transfer, a model is trained on the mixed real and
synthetic training set described in Section~\ref{sec:final-dataset-composition}
and evaluated exclusively on the real-only subset of the test split.
Table~\ref{tab:sim2real-transfer} compares this model against a variant
trained only on real data, to isolate the effect of synthetic data
augmentation on real-world performance.

\begin{table}[H]
\centering
\caption{Sim2Real transfer evaluation: performance on real-only test data
for models trained on real-only versus mixed real+synthetic data.}
\label{tab:sim2real-transfer}
\begin{tabular}{@{}lcc@{}}
\toprule
\textbf{Metric} & \textbf{Trained on Real-only} & \textbf{Trained on Mixed (Real+Synthetic)} \\
\midrule
Occupancy accuracy (\%)      & \todo{XX.X} & \todo{XX.X} \\
Classification accuracy (\%) & \todo{XX.X} & \todo{XX.X} \\
Weight MAE (kg)               & \todo{X.X}  & \todo{X.X} \\
\bottomrule
\end{tabular}
\end{table}
%% TODO: Insert actual Sim2Real transfer results and relate them back to
%% the domain gap assessment in Section~\ref{sec:domain-gap-assessment}.

\section{Discussion of Results and Failure Cases}
\label{sec:discussion-failure-cases}
Overall, the fused model is expected to outperform both single-modality
baselines across most tasks, consistent with the complementary nature of
camera and pressure information motivated in
Section~\ref{sec:motivation}. Anticipated failure cases include
ambiguous object-versus-child-seat pressure signatures under camera
occlusion, and pose estimation errors for heavily reclined seating
positions not well represented in the training data.
%% TODO: Replace with an actual discussion grounded in the final
%% experimental results, including specific qualitative failure examples.
